\section{Experimental setup}
\label{sec:experimental_setup}

\subsection{Datasets}

\paragraph{DRIVE.} The Digital Retinal Images for Vessel Extraction dataset
\citep{staal2004drive} contains 40 color fundus photographs (565 $\times$ 584
pixels) with manually annotated binary vessel masks. We use the standard
20/20 train/test split. From the 20 training images, we hold out 4 for
validation (used for threshold selection and temperature fitting).

\paragraph{STARE.} The Structured Analysis of the Retina dataset
\citep{hoover2000stare} contains 20 fundus images (700 $\times$ 605 pixels)
with vessel annotations by two experts. We use the first expert's annotations
and evaluate on all 20 images in a zero-shot cross-dataset setting (no
fine-tuning on STARE).

\paragraph{CHASE\_DB1.} The Child Heart and Health Study in England dataset
\citep{fraz2012chase} contains 28 fundus images (999 $\times$ 960 pixels)
from pediatric subjects, with vessel annotations. We again evaluate in a
zero-shot setting using all 28 images.

Using STARE and CHASE\_DB1 as external test sets tests whether uncertainty
quality, not just segmentation accuracy, transfers under domain shift. The
three datasets differ in resolution, patient population (adults vs.\ children),
imaging equipment, and annotation protocols.

\subsection{Preprocessing and augmentation}

All images are resized to 512 $\times$ 512 pixels. During training, we
extract 256 $\times$ 256 patches with vessel-aware sampling: 80\% of patches
are constrained to contain at least 16 vessel pixels, addressing the heavy
class imbalance in retinal images (vessels typically occupy 10--15\% of pixels).
We extract 32 patches per image per epoch. Training augmentations include
random horizontal and vertical flips, rotations up to 15$^\circ$, and
brightness/contrast jittering.

At test time, we process full 512 $\times$ 512 images (no patching) to
avoid boundary artifacts in the uncertainty maps.

\subsection{Model and training}

The base model is a U-Net \citep{ronneberger2015unet} with a ResNet-34
encoder \citep{he2016deep}, pretrained on ImageNet, implemented using
the \texttt{segmentation-models-pytorch} library. Dropout2d with
$p = 0.3$ is applied after the decoder and before the final $1 \times 1$
convolution.

Training uses the Adam optimizer \citep{kingma2015adam} with learning rate
$10^{-4}$, weight decay $10^{-4}$, and gradient clipping at norm 1.0.
The loss function is a hybrid of Dice loss and binary cross-entropy:
\begin{equation}
  \mathcal{L} = \frac{1}{2}\mathcal{L}_{\text{Dice}} +
  \frac{1}{2}\mathcal{L}_{\text{BCE}}
  \label{eq:loss}
\end{equation}
with a positive class weight of 4.0 in the BCE term to compensate for
class imbalance. The Dice loss uses a smoothing constant of 1.0.

We train for 80 epochs with batch size 6. The model converges around epoch
60; we select the checkpoint with the best validation Dice.

For the ensemble baseline, we train 5 U-Net models with identical
hyperparameters but different random seeds (controlling weight
initialization, data shuffling, and dropout masks).

\subsection{Evaluation protocol}

All metrics are computed on the DRIVE test set (20 images) unless otherwise
stated. For each image, we compute the mean prediction and uncertainty map
using the specified method (MC Dropout with $T = 30$ passes, or TTA with
$K = 6$ augmentations). Deferral thresholds are fit on the validation set
and applied without modification to the test set.

We report results for six configurations, covering all combinations of
\{MC Dropout, TTA\} $\times$ \{global, adaptive, confidence-aware\} deferral,
plus additional experiments with temperature scaling.

\subsection{Metrics}

\paragraph{Segmentation quality.}
\begin{align}
  \text{Dice} &= \frac{2 \sum_{ij} \hat{y}_{ij} y_{ij}}
  {2\sum_{ij} \hat{y}_{ij} y_{ij} + \sum_{ij} \hat{y}_{ij}(1 - y_{ij}) +
  \sum_{ij} (1 - \hat{y}_{ij}) y_{ij}}
  \label{eq:dice} \\
  \text{IoU} &= \frac{\sum_{ij} \hat{y}_{ij} y_{ij}}
  {\sum_{ij} \hat{y}_{ij} y_{ij} + \sum_{ij} \hat{y}_{ij}(1 - y_{ij}) +
  \sum_{ij} (1 - \hat{y}_{ij}) y_{ij}}
  \label{eq:iou} \\
  \text{AUC} &= \text{Area under the ROC curve, computed on } \hat{p}_{ij}
  \text{ vs.\ } y_{ij}
  \label{eq:auc}
\end{align}

\paragraph{Calibration.}
\begin{equation}
  \text{ECE} = \sum_{b=1}^{B} \frac{n_b}{N}
  \left| \text{acc}(b) - \text{conf}(b) \right|
  \label{eq:ece}
\end{equation}
where $B = 15$ bins, $n_b$ is the number of pixels in bin $b$,
$\text{acc}(b)$ is the fraction of correctly predicted pixels in bin $b$,
and $\text{conf}(b)$ is the mean predicted probability in bin $b$.

\paragraph{Uncertainty quality.}
\begin{equation}
  \text{Unc-AUROC} = \text{AUC}(u_{ij}, \; \mathbb{1}[\hat{y}_{ij} \neq y_{ij}])
  \label{eq:unc_auroc}
\end{equation}
This is the ROC AUC when uncertainty is used as a binary classifier of
prediction error. An Unc-AUROC of 0.5 means uncertainty is no better than
random at identifying errors; 1.0 means it perfectly separates correct from
incorrect pixels.

\paragraph{Risk-coverage metrics.} For a given deferral threshold, we compute
the Dice score on accepted pixels only and plot it against coverage (fraction
of accepted pixels). The Area Under the Coverage Curve (AUCC) summarizes this
trade-off in a single number:
\begin{equation}
  \text{AUCC} = \int_0^1 \text{Dice}(\text{coverage}) \; d(\text{coverage})
  \label{eq:aucc}
\end{equation}
Higher AUCC means the method's uncertainty is more effective at identifying
which pixels to keep.

\paragraph{Deferral metrics.} Coverage and error reduction ratio (ERR) are
defined in Section~\ref{sec:formulation} (Equation~\ref{eq:err}).
